% ==========================================================================
% GENERATED by scripts/split_paper.py. Do not edit this file.
%
% Source: imported_papers/Koopman-JEPA-TMLR-rewrite/main.tex
% Part 2: Error Growth, and What a Spectral Bound Does Not Say
%
% Edits belong upstream, in the published paper. To change what the
% chapter includes, change imports/bounded.toml and re-run
% the splitter.
% ==========================================================================
\section{Error Growth, and What a Spectral Bound Does Not Say}
\label{sec:background}

\subsection{Koopman Lifting Turns a Rollout into a Matrix Power}
\label{sec:koopman}
% [orig 2.1] DONE. Three paragraphs of the approved text compressed into two.
% Both footnotes kept, the second shortened. The Brunton citation on
% Koopman-invariant subspaces is kept because it is what makes the residual
% irreducible rather than a matter of width.
Let $\z_t \in \mathbb{R}^D$ be a latent representation of an environment at time step $t$, and let $f_\theta$ be a learned nonlinear transition model pointwise satisfying $\|f_\theta(\z) - f^*(\z)\| \le \epsilon$, where $f^*$ denotes the true latent space dynamics. To predict several steps ahead the model feeds its own output back into itself, and each application multiplies the accumulated error by the Lipschitz constant $L$ of the transition function\footnote{The Lipschitz constant $L$ of a function $f$ is the smallest value such that $\|f(\z) - f(\z')\| \le L\|\z - \z'\|$ for all inputs $\z, \z'$. It measures how much $f$ can stretch distances between inputs. If $L < 1$ the function contracts distances and small errors shrink at each step. If $L > 1$ the function stretches distances and small errors grow.}. For non-trivial physical dynamics $L > 1$ holds,\footnote{Even in globally dissipative systems the local Lipschitz constant exceeds one wherever trajectories diverge from one another, such as near an unstable equilibrium (a pendulum near its upright position) or transverse to the stable manifold of a saddle point. One such region traversed during a rollout is enough to drive exponential worst-case growth.} so the $k$-step rollout $\hat{\z}_{t+k} = f_\theta^k(\z_t)$ accumulates error exponentially:
\begin{equation}
    \|\hat{\z}_{t+k} - \z^*_{t+k}\| \;\leq\;
    \epsilon \sum_{i=0}^{k-1} L^i \;=\; \mathcal{O}(L^k \epsilon).
    \label{eq:compound}
\end{equation}
This is a property of chaining nonlinear maps rather than a training deficiency, and the bound places no limit on trajectory error as $k$ grows.

Koopman theory~\citep{Koopman1931,Mezic2005} offers a way around this compounding. Let $\x_t \in \mathcal{X}$ be the raw observation, evolving as $\x_{t+1} = F(\x_t)$. Rather than predicting the raw state, Koopman theory advances measurements taken from it. A measurement, called an \emph{observable}, is any function $g: \mathcal{X} \to \mathbb{R}$ that extracts a scalar feature from the raw state, and the Koopman operator $\mathcal{K}$ acts on it by composition:
\begin{equation}
    (\mathcal{K}\,g)(\x_t) \;=\; g(F(\x_t)) \;=\; g(\x_{t+1}).
    \label{eq:koopman_def}
\end{equation}
The operator is linear even when $F$ is not, but it acts on an infinite-dimensional space and cannot be computed directly. The practical workaround learns $D$ observables with a neural encoder $\phi$~\citep{Lusch2018,Bevanda2021}, \emph{lifting} the observation into $\z_t = [g_1(\x_t), \dots, g_D(\x_t)]^\top$, where $\mathcal{K}$ is approximated by a $D \times D$ matrix $\K$ and the transition becomes $\z_{t+1} \approx \K\z_t$. An autonomous $k$-step rollout is then one matrix power, $\hat{\z}_{t+k} = \K^k\z_t$, and with actions it stays sequential, $\z_{t+1} = \K\z_t + \B\mathbf{a}_t$. The finite lift leaves an irreducible residual $\epsilon_\text{lin}$, because systems with multiple isolated fixed points, the pendulum among them, admit no finite-dimensional Koopman-invariant subspace containing the state~\citep{Brunton2016Invariant}. The question is whether its accumulation is controlled.

\subsection{A Non-Expansive Spectrum Bounds Error Growth Linearly}
\label{sec:bound}
% [orig 2.2] DONE. Verbatim except for three changes. "Recall that" is a banned
% opener and the sentence starts on its subject instead. The emphasis on "not"
% becomes a sentence break, which carries the same weight without using \emph
% for something that is not a defined term. The forward reference points at
% Section 4.1, which is where the variants now live.
The finite-dimensional lifting leaves a single-step residual $\epsilon_\text{lin}$, so that $\|\z^*_{t+1} - \K\z^*_t\| \le \epsilon_\text{lin}$. Whether this error amplifies depends on the \emph{spectral norm} of $\K$, written $\|\K\|_2$. The spectral norm is the maximum factor by which $\K$ can stretch any non-zero vector, and is equal to the square root of the largest eigenvalue of $\mathbf{K}^\top\mathbf{K}$:
\begin{equation}
    \|\mathbf{K}\|_2 = \max_{\mathbf{z}\neq\mathbf{0}} \frac{\|\mathbf{K}\mathbf{z}\|_2}{\|\mathbf{z}\|_2} = \sqrt{\lambda_{\max}(\mathbf{K}^\top\mathbf{K})},
\end{equation}
where $\lambda_{\max}(\cdot)$ denotes the maximum eigenvalue.

If $\|\K\|_2 \le 1$, the system is stable in the sense that a small error introduced at one step cannot be magnified by further applications of $\K$. Under this condition, the $k$-step prediction error from an initial state error $\delta$ satisfies:
\begin{equation}
    \|\hat{\z}_{t+k} - \z^*_{t+k}\|
    \;\leq\; \|\mathbf{K}^k\|_2 \|\delta\| + k\,\epsilon_\text{lin}
    \;\leq\; \|\delta\| + \mathcal{O}(k\,\epsilon_\text{lin}).
    \label{eq:linear_bound}
\end{equation}

The first inequality splits the $k$-step error into the propagated initial error, amplified by at most $\|\K^k\|_2$, and the accumulated approximation error, where each of the $k$ residuals is propagated by the steps that follow it and contributes $\sum_{i=0}^{k-1}\|\K^i\|_2\,\epsilon_\text{lin}$. The second inequality applies $\|\K\|_2 \le 1$, which bounds both $\|\K^k\|_2$ by one and the sum by $k\,\epsilon_\text{lin}$. The exponential factor $L^k$ of Equation~\ref{eq:compound} is thereby replaced by the linear factor $k$.

The bound is conditional on this constraint. Linearity by itself does not bound rollout error. When $\rho(\K) > 1$, where $\rho(\K) = \max_i |\lambda_i(\K)|$ is the spectral radius, the linear iterates $\K^k$ themselves grow exponentially, so an unconstrained linear operator compounds error at rate $\rho(\K)^k$, just as a nonlinear map with Lipschitz constant $L > 1$ does. A linear model is therefore stable only when its spectrum is controlled.

Controlling $\|\K\|_2$ in practice is what motivates the variants of Section~\ref{sec:variants}: a Cayley parameterization forces $\K$ to be \emph{orthogonal}, meaning $\K^\top\K = I$, which makes $\|\K\|_2 = 1$ hold exactly, while a soft penalty constrains only the spectral radius and carries no bound of the form of Equation~\ref{eq:linear_bound}. The bound carries over to the action-conditioned setting whenever actions are known at evaluation time.\footnote{For some matrices, $\|\K\|_2$ is larger than the largest eigenvalue magnitude $\rho(\K)$. Such a matrix can stretch a vector for a few steps even though all its eigenvalues lie inside the unit circle, a phenomenon known as \emph{transient growth}~\citep{trefethen2005spectra}. Constraining $\rho(\K)$, as our soft variant does, is therefore a slightly weaker guarantee than constraining $\|\K\|_2$. For orthogonal matrices the two quantities coincide, so the strict variant is unaffected.}

\subsection{What the Bound Does Not Say}
\label{sec:bound_limits}
% [draft 2.3] DONE. The draft's prose with four style fixes: two semicolons and
% two colons become sentence breaks, and the emphasis on "grows" is replaced by
% the lexical contrast between rate and size, which carries the same point
% without using \emph for a word that is not a defined term.
Equation~\ref{eq:linear_bound} limits the rate at which error grows. It says nothing about the size of that error. The two are easy to conflate, and the distinction is what the experimental sections turn on.

The extreme case makes the gap visible. A predictor that ignores its input and always returns the same vector, $f(\z) = \mathbf{c}$, has Lipschitz constant $L = 0$, so it cannot amplify anything and its error curve is perfectly flat in $k$. It is also useless. Copying the input forward unchanged, $f(\z) = \z$, is the linear version of the same idea. It is non-expansive with $\|\K\|_2 = 1$ and satisfies the hypothesis of Equation~\ref{eq:linear_bound} exactly. Non-expansiveness is thus a property that trivial predictors have in abundance, and finding it in a learned operator is not on its own evidence that the operator has learned anything.

What the bound leaves open is the size of the residual $\epsilon_\text{lin}$. Equation~\ref{eq:linear_bound} promises only that error stays within $\|\delta\| + k\,\epsilon_\text{lin}$, which is a strong statement when $\epsilon_\text{lin}$ is small and an empty one when it is not. Constraining the spectrum does not make $\epsilon_\text{lin}$ small, and can enlarge it. Restricting $\K$ to rotations removes the freedom to contract along directions in which the true dynamics dissipate energy, so the operator that is easiest to bound need not be the one that fits best. The bound and the fit pull in opposite directions, and the theory by itself cannot say where the balance falls.

The consequence for evaluation is concrete. A flat error curve is equally consistent with a model that has captured the dynamics and with one that has stopped responding to its input, so the two must be separated by measurement rather than assumed apart. We therefore carry the two trivial predictors above through the experiments as controls, define them precisely in Section~\ref{sec:measurements}, and report them beside the learned models throughout Section~\ref{sec:exp}, most directly in Section~\ref{sec:inversion}.

% =============================================================================
